% ============================================================================
%  REVISED TEXT SNIPPETS  (copy-paste, not \input)
%
%  The current main.tex Abstract and Conclusion claim the method
%  "significantly improves" feedback. The data does NOT support that
%  (paired t-test p=0.42, Wilcoxon p=0.70, Cohen's d=0.057). Below are
%  honest replacements. Copy the body of each into the matching block of
%  main.tex. These are plain text snippets, not meant to be \input.
% ============================================================================


% ---------------------------------------------------------------------------
% (1) REPLACEMENT ABSTRACT  -- replace the text inside \begin{abstract}...\end{abstract}
% ---------------------------------------------------------------------------
%
% Preparing for software engineering interviews is a critical but
% resource-intensive process for computer science students. While large
% language models (LLMs) offer a scalable solution for automated coaching,
% their feedback often lacks the specificity and actionability required for
% high-stakes technical interviews. In this project, we build an AI Interview
% Coach based on Qwen2.5-3B-Instruct and attempt to improve it through
% Reinforcement Learning from AI Feedback (RLAIF) and Direct Preference
% Optimization (DPO), using a stronger judge model (Qwen3.5-9B) to generate
% preference data over a five-dimension coaching rubric. On a fixed 200-example
% test set, the SFT+DPO model attains a marginally higher mean judge score than
% the baseline (17.03 vs. 16.86 on a 1--20 scale), but the difference is not
% statistically significant (paired $t$-test $p=0.42$; Wilcoxon $p=0.70$;
% Cohen's $d=0.057$). A per-item failure analysis shows that genuine gains---
% notably the removal of baseline hallucinations---are offset by newly
% introduced technical errors, and that a strong ceiling effect (65\% of
% baseline outputs already score $\geq 18/20$) together with a coarse,
% homogeneous evaluation set limits the measurable benefit. We discuss the
% methodological implications for evaluating preference optimization with LLM
% judges.


% ---------------------------------------------------------------------------
% (2) REPLACEMENT CONCLUSION  -- replace the body of \section{Conclusion}
% ---------------------------------------------------------------------------
%
% We built an AI Interview Coach by fine-tuning Qwen2.5-3B with SFT and DPO on
% RLAIF-labeled preference data, and evaluated it against the baseline with an
% LLM judge on a fixed 200-example test set. The trained model scored
% marginally higher on every rubric dimension, but no difference was
% statistically significant, and the effect size was negligible. Our failure
% analysis explains why: the judge saturates near the top of its scale (a
% ceiling effect), its scores are low-resolution, the evaluation set is
% homogeneous (all test answers share a single answer-type label), and real
% improvements are cancelled by occasional new technical errors introduced
% during preference optimization. Rather than a clear win, our main
% contributions are negative but instructive: (i) LLM judges that saturate are
% weak instruments for detecting incremental alignment gains, motivating
% pairwise/win-rate evaluation and finer rubrics; (ii) reference-free DPO can
% introduce fluent but false technical claims, motivating a factuality-grounded
% preference signal. Future work includes rebuilding a diverse test set across
% all student-answer types and difficulties, a small-scale human evaluation to
% validate the judge, and grounding the coach against reference answers.


% ---------------------------------------------------------------------------
% (3) METHODOLOGY FIX  -- factual corrections to \section{Methodology}
% ---------------------------------------------------------------------------
%  - The dataset size in the text says "4,853" in the abstract/intro and the
%    methodology says "4,853"; the released split is ~4,653 train + 200 test.
%    Use the number that matches your final data (the repo README states 4,653
%    train / 200 test).
%  - The methodology lists three student-answer types as if all are present in
%    the data. In the released test/train files every item is labeled
%    "partially_correct". Either regenerate the data with the full set of types
%    or soften the claim to: "We designed prompts for five answer types; the
%    evaluation in this paper uses the partially-correct subset (see
%    Limitations)."
